%==============================================================================
% Sjabloon poster graduaatsproject
%==============================================================================
% Gegevens staan in ../projectgegevens.tex.
% De poster is een affiche van A0-formaat: hij wordt van op twee meter afstand
% bekeken, op de verdedigingsdag en op infodagen. Schrijf dus kort, gebruik
% beeld, en zorg dat iemand in dertig seconden begrijpt wat je gemaakt hebt.

\documentclass[a0,portrait]{hogent-poster}

\input{../projectgegevens}

% Info over de opleiding
\course{Graduaatsproject}
\studyprogramme{Graduaat Programmeren}
\academicyear{2026-2027 1e semester}
\institution{Hogeschool Gent, Valentin Vaerwyckweg 1, 9000 Gent}

% Info over het project
\title{\projecttitel}
\subtitle{\projectondertitel}
\author{Sarah Turner Burkitt}
\email{sarah.turnerburkitt@hogent.student.be}
\supervisor{Isabelle Demeulenaere}
\cosupervisor{Bart D'Hoore (Endare-specifiek: Vesalius.ai(health))}

% Deze informatie komt onderaan de samenvatting te staan
\specialisation{Programmeren}
\keywords{WYSIWYG e-maileditor,POC,Graduaatsproject}
\projectrepo{https://github.com/STB95/Graduaatsproject_TurnerBurkittSarah_26-27}

\begin{document}

\maketitle

\begin{abstract}
  % Drie tot vijf zinnen: voor welk bedrijf, welk probleem, wat heb je gemaakt
  % en wat levert het op. Dit is het enige stuk lopende tekst op je poster dat
  % iemand echt van begin tot eind leest.
  Vervang deze tekst door de samenvatting van je project.
\end{abstract}

\begin{multicols}{2}

\section{Probleem}

% Wat liep er mis of ontbrak er, en voor wie? Twee tot vier korte zinnen, of
% drie bullets. Geen achtergrondverhaal.

\section{Aanpak en technologie}

% Hoe heb je het aangepakt en waarmee? Noem de technologie bij naam en zeg in
% één zin waarom het die geworden is. Een klein architectuurschema doet hier
% meer werk dan een paragraaf tekst.

\begin{center}
  \captionsetup{type=figure}
  \includegraphics[width=1.0\linewidth]{voorbeeldafbeelding}
  \captionof{figure}{Vervang deze afbeelding door je architectuurschema of
    een schermafbeelding van je oplossing. Op een poster is beeld belangrijker
    dan tekst: reken op minstens twee afbeeldingen.}
\end{center}

\section{Resultaat}

% Wat is er opgeleverd, en werkt het? Toon het liefst met een schermafbeelding
% of een cijfer: "verwerkt 1200 orders per uur", "bespaart het supportteam
% ongeveer een uur per dag". Concreet verslaat algemeen.

\section{Wat ik eruit meeneem}

% Twee tot drie zinnen: wat heb je geleerd, technisch en over jezelf als
% beginnende professional? Dit is het menselijke stuk van je poster, en het is
% wat bezoekers op een infodag onthouden.

\end{multicols}
\end{document}
